\section{Timeline of substantive modifications}

This timeline selects changes that altered the calendar's architecture, precedence, or widely observed contents. It is not a list of every canonization or feast inserted between editions.

\begin{longtable}{@{}>{\RaggedRight\arraybackslash}p{0.13\linewidth}>{\RaggedRight\arraybackslash}p{0.27\linewidth}>{\RaggedRight\arraybackslash}p{0.53\linewidth}@{}}
\toprule
\textbf{Date} & \textbf{Act or book} & \textbf{Substantive effect}\\
\midrule
\endfirsthead
\toprule
\textbf{Date} & \textbf{Act or book} & \textbf{Substantive effect}\\
\midrule
\endhead
\bottomrule
\endfoot
1568 & Pius V, \emph{Quod a nobis}; reformed Roman Breviary & Revised and standardized the calendar as embedded in the Roman Office, while defining the books and uses to which the reform applied.\\
1570 & Pius V, \emph{Quo primum}; \emph{Missale Romanum} & Published the corresponding Roman Missal and its annual Temporale--Sanctorale structure for broad use under the bull's stated conditions.\\
1604 and 1634 & Typical Missal editions under Clement VIII and Urban VIII & Revised the post-Tridentine Missal and continued the authoritative editorial development of its calendar and rubrics.\\
1911--1913 & Pius X, \emph{Divino afflatu}; new rubrics and calendar adjustments & Reordered the weekly psalter and precedence, curtailed the habitual displacement of ferial offices, and coordinated later calendar changes with the new system.\\
1951 and 1955 & Experimental Easter Vigil and the decree \emph{Maxima redemptionis nostrae mysteria} & Restored the rites of Holy Week to revised times and forms; the Sacred Triduum and Paschal Vigil again became a distinctive liturgical unit in practice.\\
1955--1956 & \emph{Cum hac nostra aetate} and the simplified rubrics & Suppressed most vigils and octaves, simplified ranks and commemorations, and substantially reduced the calendar before the final 1960 codification.\\
25 July 1960; effective 1 January 1961 & John XXIII, \emph{Rubricarum instructum} and \emph{Codex rubricarum} & Established four classes, new precedence and occurrence rules, and a coordinated calendar for the Roman Breviary and Missal. This is the immediate normative system of the 1962 book.\\
1961 & Sacred Congregation of Rites, instruction on particular calendars & Required diocesan and religious calendars to be revised under the new code, distinguishing the general calendar from properly local claims.\\
23 June 1962 & Typical edition of the \emph{Missale Romanum} & Published the Missal whose printed calendar and propers form this study's edition baseline.\\
14 February 1969; effective 1 January 1970 & Paul VI, \emph{Mysterii Paschalis}, and the revised General Roman Calendar & Reconstructed the general calendar for the postconciliar Roman Rite under new universal norms. It is a successor calendar, not an update to be merged into the 1962 tables.\\
19 March 2020 & Congregation for the Doctrine of the Faith, \emph{Cum sanctissima} & For celebrations according to the 1962 books, permitted defined optional celebrations of later-canonized saints and certain entries in the \emph{pro aliquibus locis} supplement. It left the printed 1962 calendar intact and imposed protected-day limits.\\
\end{longtable}

The 1969--1970 reform is included to mark the historical boundary. Its nomenclature and options must not be projected backward: a postconciliar ``optional memorial'' is not the same thing as a 1962 commemoration, III-class feast, or later permission under \emph{Cum sanctissima}.
